% ===================================================================
%  図表第 2 号 — 人体。— 休み時間。— 遊び。
%
%  Ceci n'est PAS une transcription : c'est une traduction, faite en
%  2026, en japonais d'aujourd'hui. Voir texto/ja/10-zuhyo-01.tex
%  pour la consigne, qui vaut pour les seize fichiers.
% ===================================================================

%%K t02-apar-1 apar
\VUsaut{4.0mm}
\VUtitre{15.0pt}{60}{図表第 2 号}
\VUsaut{2.0mm}
\VUtitre{11.4pt}{0}{人体。--- 休み時間。}
\VUtitre{11.4pt}{0}{遊び。}

%%K t02-tit-0 sub
\VUtitre{13.2pt}{0}{人体。}
\VUsaut{2.4mm}
\VUcentre{10.2pt}{0}{\textit{（やさしい博物の一課。）}}
\VUsaut{3.0mm}

%%K t02-01-1 p
1. --- 人体の主な部分は次のとおりである。\VUgras{頭}、\VUgras{胴}、
\VUgras{四肢}。

%%K t02-tit-1 sub
\VUpk{10.2pt}{40}{頭。}
\VUsaut{3.0mm}

%%K t02-02-1 p
2. --- 頭は二つの部分に分かれる。\VUgras{頭蓋}、その中には\VUgras{脳}が
あり、外は\VUgras{髪}が覆っている。それに\VUgras{顔}である。

%%K t02-03-1 p
3. --- 我々の絵では頭蓋も脳も見えない。しかし顔の大切な部分はたいそうはっきり
見える。

%%K t02-04-1 p
4. --- まず\VUgras{額}がある。強い才知と、絶えず働く心とを示している。
\VUgras{こめかみ}はたいそう広い。\VUgras{目}は生き生きとして大きく開いている。
濃い\VUgras{眉}と長い\VUgras{まつげ}がそれを守っている。

%%K t02-05-1 p
5. --- ここに\VUgras{鼻}と\VUgras{鼻孔}もある。鼻は匂いを嗅ぎ、息をするのに
役立つ。鼻の両側には\VUgras{頬}があり、その先には\VUgras{耳}がある。顔の下の
部分は\VUgras{口}と\VUgras{顎}とでできている。

%%K t02-06-1 p
6. --- この二つはよくは見えない。\VUgras{口ひげ}と\VUgras{あごひげ}が隠して
いるからである。しかし我々は知っている。口には二枚の\VUgras{唇}、
\VUgras{上顎}と\VUgras{下顎}がその\VUgras{歯}とともにあり、\VUgras{舌}がある
ことを。口で我々は話し、食べる。舌で我々は食べ物の味を知る。

%%K t02-07-1 p
7. --- 目、耳、鼻、口は、指とともに、五つの感覚の器官である。\VUgras{視覚}、
\VUgras{聴覚}、\VUgras{嗅覚}、\VUgras{味覚}、\VUgras{触覚}である。

%%K t02-08-1 p
8. --- こうして頭は人体のうちで最も面白い部分の一つである。

%%K t02-tit-2 sub
\VUsaut{4.4mm}
\VUpk{10.2pt}{40}{胴。}
\VUsaut{3.0mm}

%%K t02-09-1 p
9. --- \VUgras{首}が頭を胴につないでいる。そこには\VUgras{のど}と
\VUgras{うなじ}が入る。

%%K t02-10-1 p
10. --- 胴にも大切な器官がたくさん入っている。\VUgras{胸}には\VUgras{肺}が
あり、それで我々は息をする。それに\VUgras{心臓}があり、それが命の中心である。
心臓が打つのをやめると、生き物は死ぬ。

%%K t02-11-1 p
11. --- \VUgras{肋骨}が胸を支えている。

%%K t02-12-1 p
12. --- 胴のほかの部分は\VUgras{脇腹}、\VUgras{背}、\VUgras{肩}、\VUgras{腹}
（お腹）である。眠るには横向きか仰向けに寝、荷は肩で担ぐ。

%%K t02-13-1 p
13. --- 腹には\VUgras{胃}と\VUgras{腸}がある。それらは食べ物を消化するのに
役立つ。

%%K t02-tit-3 sub
\VUsaut{4.4mm}
\VUpk{10.2pt}{40}{四肢。}
\VUsaut{3.0mm}

%%K t02-14-1 p
14. --- 我々には四つの\VUgras{肢}がある。二本の\VUgras{腕}と二本の
\VUgras{脚}である。腕で我々は物を取り、握り、持ち上げ、拾う。脚で我々は立ち、
歩き、走る。

%%K t02-15-1 p
15. --- \VUgras{肩}が腕を体につないでいる。たいそう丈夫な\VUgras{筋肉}、
すなわち二頭筋が腕を動かし、\VUgras{肘}が腕を\VUgras{前腕}につなぐ。
\VUgras{手首}が\VUgras{手}の関節をなし、手の先は\VUgras{指}と\VUgras{爪}で
終わる。手には\VUgras{静脈}が一本（それに動脈が一本）見え、そこを\VUgras{血}
が流れている。

%%K t02-16-1 p
16. --- 脚の部分は次のとおりである。\VUgras{もも}、\VUgras{ひざ}と
\VUgras{ひかがみ}、\VUgras{ふくらはぎ}、その\VUgras{肉}を\VUgras{皮}が覆って
いる。それに\VUgras{くるぶし}と\VUgras{足}である。脚の\VUgras{骨}はたいそう
太く、たいそう強い。足の先には\VUgras{足の指}がある。ほかの部分は
\VUgras{足の裏}と\VUgras{かかと}である。

%%K t02-17-1 p
17. --- 以上が人体のほぼ完全な説明である。

%%K t02-17-2 p
\textit{注。} --- \textit{「私の……が痛い（頭、など）」という言い方を借りて、
先生は} \VUgras{体の具合}\textit{のよしあしという見方から、この語彙をひととおり
おさらいできるであろう。}

%%K t02-tit-4 sub
\VUsaut{5.0mm}
\VUpk{10.2pt}{40}{体操。}
\VUsaut{2.4mm}
\VUcentre{10.2pt}{0}{\textit{（生徒の一人が語る。）}}
\VUsaut{3.0mm}

%%K t02-18-1 p
18. --- 体を柔らかく、身軽に、丈夫に保つには、脚と腕を働かせねばならない。
だから我々は遊び、また\VUgras{体操}をする。

%%K t02-19-1 p
19. --- ふだんこの二つは学校の中庭でする（図表第 1 号を見よ）。しかし木曜と
日曜には広い草地へ行く。そこには走る場所も、楽しむ場所もある。

%%K t02-20-1 p
20. --- 先生や父母がときには一緒に来る。しかしこの練習を指揮するのは
\VUgras{体操の先生} \textsuperscript{(12)} である。

%%K t02-21-1 p
21. --- 草地の入口の左には\VUgras{体操器械} \textsuperscript{(1)} が立っている。
我々は\VUgras{梯子} \textsuperscript{(2)} を登り、\VUgras{吊り輪}
\textsuperscript{(3)} や\VUgras{ぶらんこ棒} \textsuperscript{(4)} で逆さに揺れ、
\VUgras{縄梯子} \textsuperscript{(5)} や\VUgras{結び目のある綱} \textsuperscript{(6)}
の先まで手だけで登る。

%%K t02-22-1 p
22. --- そのあいだ仲間たちは\VUgras{木馬} \textsuperscript{(7)} や
\VUgras{平行棒} \textsuperscript{(11)} を跳び越える。ほかの者は\VUgras{拳闘}
\textsuperscript{(8)} をし、あるいは面をつけ\VUgras{フルーレ}を持って
\VUgras{剣術} \textsuperscript{(9)} をする。また\VUgras{亜鈴} \textsuperscript{(10)}
を上げる者もいる。

%%K t02-tit-5 sub
\VUsaut{4.6mm}
\VUpk{10.2pt}{40}{遊び。}
\VUsaut{3.0mm}

%%K t02-23-1 p
23. --- 体操をする者のまわりでは、男の子女の子がいろいろな\VUgras{遊び}を
している。女の子は自分の\VUgras{輪} \textsuperscript{(14)} で遊ぶ。
\VUgras{縄跳び} \textsuperscript{(13)} をし、あるいは兄弟や従姉妹を
\VUgras{クロケー} \textsuperscript{(15)} の一勝負に誘う。相手が門を楽々くぐれる
と思ったちょうどそのとき、彼女たちは\VUgras{木槌}で相手の\VUgras{球}を打ち
のけるのだから、それは愉快でならない。

%%K t02-24-1 p
24. --- 男の子はもっと力のいる遊びが好きである。喜んで\VUgras{九柱戯}
\textsuperscript{(16)} をし、\VUgras{球} \textsuperscript{(17)} で手際よく柱を倒す。
\VUgras{独楽} \textsuperscript{(18)} や\VUgras{剣玉} \textsuperscript{(19)} も、
\VUgras{蛙投げ} \textsuperscript{(20)} さえも軽んじはしない。しかし一番好きな
遊びは\VUgras{おはじき} \textsuperscript{(21)} と\VUgras{手球} \textsuperscript{(22)}
である。\VUgras{温室} \textsuperscript{(26)} の壁がちょうど軽い球を弾ませるのに
よいからである。

%%K t02-25-1 p
25. --- 小さなジャンは自分の\VUgras{竹馬} \textsuperscript{(24)} に乗ってたいそう
巧みで、うまく釣合を保っている。そのそばでは何人かの仲間がその温室の中や
\VUgras{生垣}の後ろで\VUgras{かくれんぼ} \textsuperscript{(25)} をしている。
\VUgras{手の甲当て} \textsuperscript{(28)} や\VUgras{羽根} \textsuperscript{(29)} の
方が好きな者もいる。羽根は一つの\VUgras{羽子板}からもう一つの羽子板へと飛んで
ゆく。

%%K t02-noto-1 noto
\VUnotes{143.2mm}{%
(*) 子供の遊びの名は、たいてい国ごとのものでしかない。我々はできるかぎり
イド語で名をつけた。あるときはその遊びを最もよく表す動作から、あるときは
誤解を招かぬ限りでどこかの国の名を取って。たとえば、\VUgras{樽の遊び}
（ドイツとフランスの）は\VUgras{蛙投げ} \textsuperscript{(20)} （イギリスの）で
あって、遊ぶ者は穴のあいた箱（樽）の上に口を開けて坐る金属の蛙の口に丸い札を
投げ入れようとする。\VUgras{肩跳び} \textsuperscript{(30)} が\VUgras{馬跳び}と
違うのは一点だけである。頭の上を跳び越えるとき、遊ぶ者は両手を仲間の肩に
つく。「\VUgras{馬跳び}」ではその背にしかつかない。--- あるいはまた、何人かが
身をかがめ、ほかの者が両手をその肩について背を跳び越える。前の者は崩れずに
衝撃を受け止めねばならず、後の者は釣合を失わずに跳ばねばならない（掛図その
ものを見よ）。%
}

%%K t02-26-1 p
26. --- 草地の真中に木が二本ある。その一本の下で七八人の男の子が
\VUgras{肩跳び} \textsuperscript{(30)} (*) を始めた。この遊びはどの国にもあるわけ
ではない。しかし\VUgras{馬跳び}が何であるかは誰でも知っている。子供たちは今は
それをしていない。

%%K t02-tit-6 sub
\VUsaut{5.0mm}
\VUpk{10.2pt}{40}{野外の遊び。}
\VUsaut{3.4mm}

%%K t02-27-1 p
27. --- \VUgras{目隠し鬼} \textsuperscript{(31)} もたいそう面白い。女の子男の子が
かわるがわる目に\VUgras{目隠し}を当て、捕まってしまう不器用な者を捕える。

%%K t02-28-1 p
28. --- 草地の真中に、いかにもフランスらしい遊びがある。少々荒っぽくはあるが。
それは\VUgras{熊} \textsuperscript{(32)} であって、あの静かな\VUgras{凧}
\textsuperscript{(33)} の遊びよりも、イギリスの\VUgras{庭球} \textsuperscript{(34)}
よりも、はるかに騒がしい。

%%K t02-29-1 p
29. --- この最後の運動は年上の生徒の楽しみである。我々の先生方も自ら喜んで
なさる。たいそうな腕前と身軽さが要り、私はそれがとても好きである。
\VUgras{ぶらんこ} \textsuperscript{(35)} は女の子に譲っておく。

%%K t02-30-1 p
30. --- もう一つのイギリスの遊びも私は好きでない。それは危なくなりかねない。

%%K t02-30-1b p
私が言うのは\VUgras{蹴球} \textsuperscript{(36)} のことで、それは兄の楽しみで
ある。私は我々の古い遊び\VUgras{陣取り} \textsuperscript{(38)} の方が好きだ。同じ
だけ体によく、はるかに手荒でない。

%%K t02-tit-7 sub
\VUsaut{4.6mm}
\VUpk{10.2pt}{40}{自転車と球の遊び。}
\VUsaut{3.0mm}

%%K t02-31-1 p
31. --- \VUgras{自転車} \textsuperscript{(37)} で草地を一巡りするのも私は好きで
ある。

%%K t02-32-1 p
32. --- 来年は\VUgras{乗馬} \textsuperscript{(39)} を習うつもりだ。馬が姿よく歩み、
また速歩で進み、駈足で障害を跳び越えるさまは、なんと美しいことか。

%%K t02-33-1 p
33. --- 夏には\VUgras{泳ぎ} \textsuperscript{(40)} を習う。川の澄んだ冷たい水に
心ゆくまで飛び込み、水浴びをする。ときには終わりに紙の\VUgras{風船}
\textsuperscript{(41)} を放つ。熱い空気で膨れるやいなや、それは厳かに空へ昇って
ゆく。

%%K t02-34-1 p
34. --- ほかにも遊びはたくさんあるが、数え上げていてはきりがない。この短い
話からも分かるとおり、我々のあの美しい草地では木曜が退屈とはほど遠いのである。

%%K t02-34-2 p
N.~B. --- \textit{先生はこの章をたやすく補うことができよう。大切な遊びを一つ
ずつもっと細かく説明すればよい。}
\VUsaut{6.0mm}
\VUfilet{22mm}
